% =============================================================================
%  Preambulo comun - Formato APA 7 (estudiante)
%  Times New Roman 12 pt, texto negro, margenes de 1 pulgada, interlineado
%  doble, sangria de primera linea de media pulgada y numeracion superior
%  derecha, conforme a APA (2020, secciones 2.19 a 2.24).
% =============================================================================

\usepackage[utf8]{inputenc}
\usepackage[T1]{fontenc}
\usepackage[spanish,es-nodecimaldot,es-tabla]{babel}
\usepackage{mathptmx}          % Times New Roman (texto y matematicas)
\usepackage[letterpaper,margin=1in]{geometry}
\usepackage{setspace}
\usepackage{graphicx}
\usepackage{array}
\usepackage{amsmath}
\usepackage{titlesec}
\usepackage{fancyhdr}
\usepackage{enumitem}
\usepackage{xcolor}
\usepackage{float}
\usepackage[hidelinks]{hyperref}

% --- Color del texto: negro en todo el documento ---------------------------
\color{black}
\definecolor{negro}{gray}{0}

% --- Alineacion --------------------------------------------------------------
% El documento sale JUSTIFICADO, que es la convencion habitual en los trabajos
% academicos en espanol. APA 7 (seccion 2.24) prescribe, en rigor, alineacion a
% la izquierda con margen derecho irregular: para adoptarla basta con quitar el
% comentario de las dos lineas siguientes.
% \usepackage[document]{ragged2e}
% \setlength{\RaggedRightParindent}{0.5in}

% --- Interlineado y sangria --------------------------------------------------
% Interlineado doble y sangria de primera linea de media pulgada en todo el
% cuerpo, salvo en las tablas, las figuras y sus notas.
\doublespacing
\setlength{\parindent}{0.5in}
\setlength{\parskip}{0pt}
% Holgura extra para que el texto justificado no se salga del margen
\setlength{\emergencystretch}{1.5em}
% Particion correcta de palabras inglesas de las referencias
\hyphenation{mis-in-ter-pre-ta-tions ho-mo-sce-das-tic-i-ty}

% --- Numeracion de pagina: superior derecha ---------------------------------
\setlength{\headheight}{14.5pt}
\pagestyle{fancy}
\fancyhf{}
\fancyhead[R]{\normalfont\normalsize\thepage}
\renewcommand{\headrulewidth}{0pt}

% --- Niveles de titulo APA 7 -------------------------------------------------
% Nivel 1: centrado, negrita.  Nivel 2: izquierda, negrita.
% Nivel 3: izquierda, negrita cursiva.
\titleformat{\section}{\normalsize\bfseries\centering}{}{0pt}{}
\titleformat{\subsection}{\normalsize\bfseries\raggedright}{}{0pt}{}
\titleformat{\subsubsection}{\normalsize\bfseries\itshape\raggedright}{}{0pt}{}
\titlespacing*{\section}{0pt}{0pt}{0pt}
\titlespacing*{\subsection}{0pt}{0pt}{0pt}
\titlespacing*{\subsubsection}{0pt}{0pt}{0pt}
\setcounter{secnumdepth}{0}

% --- Encabezados de tabla y figura al estilo APA -----------------------------
% Numero en negrita, en linea propia; titulo en cursiva, en la linea siguiente.
\newcommand{\apatabla}[2]{%
  \noindent{\normalsize\textbf{Tabla #1}}\par
  \noindent{\normalsize\textit{#2}}\par
  \vspace{4pt}}

\newcommand{\apafigura}[2]{%
  \noindent{\normalsize\textbf{Figura #1}}\par
  \noindent{\normalsize\textit{#2}}\par
  \vspace{4pt}}

% Nota de tabla o figura, a espacio sencillo. Se usa un grupo con
% \setstretch{1} y no el entorno spacing de setspace: ese entorno, al cerrarse
% dentro de un documento a doble espacio, inserta un salto negativo que monta
% el texto siguiente sobre la tabla.
\newcommand{\apanota}[1]{%
  \par\vspace{4pt}%
  \begingroup\setstretch{1}\footnotesize
  \noindent\textit{Nota.} #1\par
  \endgroup}

% --- Entorno de referencias con sangria francesa -----------------------------
\newenvironment{referencias}
  {\begin{list}{}{%
     \setlength{\leftmargin}{0.5in}
     \setlength{\itemindent}{-0.5in}
     \setlength{\listparindent}{0pt}
     \setlength{\labelwidth}{0pt}
     \setlength{\labelsep}{0pt}
     \setlength{\itemsep}{0pt}
     \setlength{\parsep}{0pt}
     \setlength{\topsep}{0pt}}}
  {\end{list}}

% --- Espaciado sencillo dentro de tablas y figuras ---------------------------
\newenvironment{apacuerpo}{\par\begingroup\setstretch{1}}{\par\endgroup}

% --- Bloque indivisible: encabezado, tabla o figura y nota en la misma pagina
\newenvironment{bloque}{\par\noindent\begin{minipage}{\linewidth}}%
  {\end{minipage}\par}
